\documentclass[12pt]{article}
\usepackage[margin=0.75in]{geometry}
\usepackage{fontspec}
\setmainfont{Latin Modern Roman}
\usepackage{microtype}
\usepackage{multicol}
\usepackage{fancyhdr}
\usepackage{titlesec}
\usepackage{enumitem}
\usepackage{xcolor}
\usepackage[hidelinks]{hyperref}
\setlength{\columnsep}{0.28in}
\setlength{\parindent}{1.1em}
\setlength{\parskip}{0.35em}
\setlength{\headheight}{15pt}
\titleformat{\section}{\large\bfseries\scshape}{}{0pt}{}
\titleformat{\subsection}{\normalsize\bfseries}{}{0pt}{}
\pagestyle{fancy}
\fancyhf{}
\fancyfoot[C]{\thepage}
\renewcommand{\headrulewidth}{0.4pt}
\newcommand{\focusbox}[1]{\par\vspace{0.4em}\noindent\begingroup\setlength{\fboxsep}{6pt}\colorbox{gray!12}{\begin{minipage}{\dimexpr\linewidth-2\fboxsep\relax}#1\end{minipage}}\endgroup\par\vspace{0.5em}}
\newcommand{\studentline}{\vspace{0.35em}\hrule\vspace{0.65em}}

\fancyhead[L]{Whately Elements of Logic}
\fancyhead[R]{Teacher Guide}
\begin{document}
\begin{center}
{\LARGE\bfseries Teacher Guide: Week 2}\par\vspace{0.25em}
{\large\scshape One Word, Two Ideas: Why Ambiguity Breaks Arguments}\par\vspace{0.4em}
{12th Grade Long Logic Unit Based on Richard Whately's \textit{Elements of Logic}}\par\vspace{0.6em}
\rule{0.82\textwidth}{0.4pt}
\end{center}

\section*{Week 2 Overview}
Week 2 builds directly on the four-step analysis habit from Week 1 by focusing on one of the most common reasons arguments fail the “test” step: ambiguous or shifting terms. Students learn to notice when a key word is being used in two different senses (equivocation) and to fix the meaning of important terms before evaluating an argument. The week draws from Whately's treatment of terms and ambiguity while remaining practical and example-driven. In the revised twelfth-grade plan, students begin a recurring habit of testing words from \textit{Macbeth} and \textit{Julius Caesar}: fate, manhood, honor, ambition, liberty, and tyranny.

\section*{Essential Question}
How can the same word mean two different things inside a single argument, and what can we do about it?

\section*{Student-Facing Materials}
\begin{itemize}[leftmargin=*]
\item \textit{One Word, Two Ideas: Why Ambiguity Breaks Arguments} — student reading handout.
\item \textit{Week 2 Argument Lab: Fixing the Term} — practice identifying ambiguous terms, repairing arguments, and creating original examples.
\item \textit{Macbeth Lit-Anchor: Words That Lie Like Truth} — Shakespeare anchor passage for testing Macbeth's confidence in ambiguous prophecies: ``none of woman born'' and ``Birnam wood ... shall come.''
\item \textit{Week 2 Logic Test: The Case of the Two-Faced Word} — puzzle-style assessment requiring students to identify equivocation, fix terms, retest arguments, and review the Week 1 four-step argument diagnosis.
\end{itemize}

\section*{Learning Goals}
Students will be able to:
\begin{enumerate}[leftmargin=*]
\item Define a term and explain how the same word can refer to different ideas.
\item Distinguish ambiguity from equivocation.
\item Identify when an argument relies on an unstated shift in the meaning of a key term.
\item Fix the meaning of a disputed term before testing an argument.
\item Repair an equivocal argument by making the term precise.
\item Apply the same term-fixing habit in a puzzle-test setting with fresh, unfamiliar notices.
\item Recognize that some apparent disagreements are actually disagreements about the meaning of words.
\item Build a portfolio term investigation using either a contemporary argument or a recurring Shakespeare anchor.
\end{enumerate}

\section*{Key Vocabulary}
term, meaning, ambiguity, equivocation, fix the term, context

\section*{Weekly Lesson Sequence}

\subsection*{Day 1: Opening Question and First Read}
\begin{itemize}[leftmargin=*]
\item Begin with: “Can two people use the exact same word and mean completely different things?” Have students write a quick example.
\item Introduce the idea that many arguments look strong only because a word quietly changes meaning.
\item Read the first half of the student reading aloud or in pairs.
\item Stop after the “bank” example. Ask: Why does the argument sound silly once the two meanings are pointed out?
\end{itemize}

\subsection*{Day 2: Finish Reading and Annotate Examples}
\begin{itemize}[leftmargin=*]
\item Finish the reading.
\item Students mark every example and label the two different meanings of the key term.
\item Complete Reading Focus questions 1–4.
\end{itemize}

\subsection*{Day 3: Term-Fixing Practice}
\begin{itemize}[leftmargin=*]
\item Use the Argument Lab. Students work in pairs to identify the ambiguous term in each argument and propose a fixed definition.
\item Discuss: How does fixing the term change whether the conclusion follows?
\end{itemize}

\subsection*{Day 4: Macbeth Lit-Anchor: Words That Lie Like Truth}
\begin{itemize}[leftmargin=*]
\item Read the guided Act 4, Scene 1 apparitions passage from the lit-anchor.
\item Put two phrases on the board: ``none of woman born'' and ``Birnam wood ... shall come.'' For each, ask: What meaning does Macbeth assume? What alternate meaning remains possible?
\item Emphasize that Macbeth's confidence depends less on a simple falsehood than on an unfixed phrase taken in the most comforting possible sense.
\item Students begin the \textit{Macbeth Equivocation Chart} portfolio artifact.
\end{itemize}

\subsection*{Day 5: Puzzle Test and Debrief}
\begin{itemize}[leftmargin=*]
\item Administer \textit{The Case of the Two-Faced Word}. Remind students that the test is not asking them to be clever with puns; it is asking them to fix terms before trusting an argument.
\item Students should complete the puzzle story and main diagnosis independently. If time remains, debrief one notice by asking: Which word changed meaning? Which conclusion stopped following once the term was fixed?
\item Use the Macbeth bridge to connect the test back to the lit-anchor: Macbeth's error is not merely that he hears false words, but that he treats one possible meaning as the only possible meaning.
\item Students add the test or its corrected version to the logic portfolio as evidence of their term-investigation skill.
\end{itemize}

\section*{Lit-Anchor Teaching Notes}
The Week 2 lit-anchor gives students the central literary version of equivocation. Macbeth is not merely tricked by statements that are false on their face. He is ruined by phrases that keep one promise to his ear while allowing another meaning in reality. That makes the passage especially useful for Whately's lesson: fix the terms before trusting the conclusion.

\begin{itemize}[leftmargin=*]
\item \textbf{Guided passage:} \textit{Macbeth} Act 4, Scene 1, the apparitions' warnings and promises.
\item \textbf{Independent passage:} Act 5, Scene 8, Macbeth and Macduff, ending with Macbeth's recognition that the fiends ``palter with us in a double sense.''
\item \textbf{Central logical pressure:} Macbeth hears broad protection where the words allow narrow exceptions.
\item \textbf{Portfolio artifact:} Macbeth equivocation chart: phrase, Macbeth's interpretation, alternate meaning, hidden assumption, effect on reasoning, and connection to Whately.
\end{itemize}

For ``none of woman born,'' students should see Macbeth's assumed meaning: no human being can harm him. Macduff exposes the alternate meaning: not born in the ordinary manner. For ``Birnam wood ... shall come,'' Macbeth assumes a whole forest must uproot itself; the words can also be fulfilled by soldiers carrying branches. The common hidden assumption is that Macbeth's first comfortable meaning is the only possible meaning.

\section*{Notes on Examples}
The student reading deliberately uses a mix of light, obvious examples (“bank”) and more serious, debatable examples (“liberal education,” equality, and tolerance). The lab adds \textit{Macbeth} and \textit{Julius Caesar} examples so students see that ambiguity is not only a trick in silly arguments but a real issue in literature, public speech, and substantive disagreements. Encourage students to bring their own examples from current events, school debates, or Shakespeare.

\section*{Logic Test Answer Guide: The Case of the Two-Faced Word}
The test assesses whether students can fix terms before applying the Week 1 follow-test. Strong answers should name both meanings precisely, rewrite the argument so the two-faced word has one stable meaning, and then explain whether the rewritten conclusion follows.

\subsection*{Part I: Main Notices}
\begin{itemize}[leftmargin=*]
\item \textbf{Fair Judge:} The two-faced word is \textit{fair}. In the first use, \textit{fair} means just, impartial, or conducted according to proper rules; in the second use, it means pleasant, clear, or not rainy. This is equivocation because the argument moves from one meaning to another. A strong rewrite keeps one meaning steady. For the trial sense: ``The trial was supposed to be just and impartial. But the judge refused to let one side present evidence. Therefore the trial was not just and impartial.'' For the weather sense: ``The weather was supposed to be fair. But the sky was not clear that day, for it rained. Therefore the weather was not fair.'' The first rewrite is stronger for judging the trial; the second is logically clearer but changes the subject to weather.
\item \textbf{Free Hall:} The two-faced word is \textit{free}. First it means liberty or protected permission to speak; then it means without monetary cost. A strong rewrite must keep one sense: ``All students should have protected speech in the Hall. Therefore no student should be punished for stating an unpopular view in the Hall.'' If a student keeps the cost sense, the argument must become about price rather than liberty.
\item \textbf{True Heir:} The two-faced word is \textit{true}. First it means rightful or genuine; then it means accurate or honest. A strong rewrite keeps one sense: ``Only the rightful heir may open the cabinet. Nora has the signed inheritance record naming her the rightful heir. Therefore Nora may open the cabinet.'' Merely saying that Nora spoke accurately does not prove rightful heirship.
\end{itemize}

\subsection*{Part II: Strengthen One Argument}
Accept any strengthened version that adds a real fact, observation, or testimony while keeping the key term in one stable sense. Students should see that fixing a word prevents equivocation, but an argument may still need evidence. For the Fair Judge notice, useful added evidence would concern the conduct of the trial, not the weather: one side was silenced, the judge ignored evidence, or the rules were applied unevenly. For the Free Hall notice, evidence should match either protected speech or cost. For the True Heir notice, evidence should establish rightful heirship, not mere honesty.

\subsection*{Part III: Week 1 Review}
For the gold-ink claim, the conclusion is that the rule must be wise. The stated reason is that the Head Tutor wrote it in gold ink. The hidden assumption is that an impressive appearance or official presentation proves wisdom. The conclusion does not follow; gold ink may draw attention, but it does not establish wisdom.

\subsection*{Part IV: Logic's Limits}
Logic can test whether \textit{fair} is being used in two senses, whether a conclusion follows after \textit{free} is fixed, and whether a hidden assumption is needed. Logic alone cannot decide whether the Head Tutor was morally wise or whether Nora actually has a legal right to the cabinet without further facts or testimony.

\subsection*{Part V: Macbeth Bridge}
Macbeth assumes that ``none of woman born'' means no human being can harm him. The phrase also allows the narrower meaning: someone not born in the ordinary manner. Macbeth's hidden assumption is that his comforting interpretation is the only possible interpretation. This is dangerous equivocation because the phrase is not simply meaningless or openly false; it works by carrying a double sense that Macbeth fails to fix.

\section*{Assessment}
Use the Week 2 logic test as a portfolio assessment of term investigation. Strong work should:
\begin{itemize}[leftmargin=*]
\item identify the key ambiguous term;
\item state both meanings accurately;
\item explain why equivocation weakens the argument;
\item rewrite the argument so the two-faced word has one clear meaning;
\item strengthen one argument by adding a real fact, observation, or testimony;
\item preserve the Week 1 distinction among conclusion, stated reason, hidden assumption, and follow-test;
\item distinguish what logic can test from what requires further facts, testimony, or judgment.
\end{itemize}

\section*{Connection to Future Weeks}
This week prepares students for later work on propositions, syllogisms, fallacies (many of which rely on equivocation), and verbal versus real disputes. The habit of fixing terms will also make the study of propositions and syllogisms clearer when those weeks arrive.
\end{document}